\subsection{SUBSETS STABLE UNDER AN ACTION. INDUCED ACTION}

DEFINITION 2. A subset \(\mathbf{A}\) of a set \(\mathbf{E}\) is called stable under an action \(\alpha \mapsto f_{\alpha}\) of \(\Omega\) on \(\mathbf{E}\) if \(f_{\alpha}(\mathbf{A}) \subset \mathbf{A}\) for all \(\alpha \in \Omega\) . An element \(x\) of \(\mathbf{E}\) is called invariant under an element \(\alpha\) of \(\Omega\) if \(f_{\alpha}(x) = x\) .

The intersection of a family of stable subsets of E under a given action is stable. There therefore exists a smallest stable subset of E containing a given subset X of E; it is said to be generated by X; it consists of the elements \((f_{\alpha_{1}} \circ f_{\alpha_{2}} \circ \cdots \circ f_{\alpha_{n}})(x)\) , where \(x \in X, n \geqslant 0, \alpha_{1} \in \Omega\) for all i.

Remark. Let E be a magma with law denoted by \(\top\) . It should be noted that a subset A of E which is stable under the left action on E on itself is not necessarily stable under the right action of E on itself; a subset A of E stable under the left (resp. right) action of E on itself is stable under the law on E but the converse is not in general true. More precisely, A is stable under the law on E if and only if \(A \top A \subset A\) whereas A is stable under the left (resp. right) action on E on itself if and only if \(E \top A \subset A\) (resp. \(A \top E \subset A\) ).

Example. Take the magma E to be the set N with multiplication. The set \(\{1\}\) is stable under the internal law of N, but the stable subset under the action of N on itself generated by \(\{1\}\) is the whole of N.

DEFINITION 3. Let \(\alpha \mapsto f_{\alpha}\) be an action of \(\Omega\) on E and A a stable subset of E. The mapping which associates with an element \(\alpha \in \Omega\) the restriction of \(f_{\alpha}\) to A (considered as a mapping of A into itself) is an action of \(\Omega\) on A said to be induced by the given action.

